\section{Supplementary Materials}
\subsection{Supplementary Methods}

\subsubsection{Deriving population-level quantities from individual-level components}

\subsubsection{Simulation approach}

We used stochastic, individual-based epidemic simulations to investigate the impact of individual-level infectiousness timing on epidemic dynamics. The simulation model has a modular design: an infection attempt generator produces a set of transmission attempts for each infected individual, and a population-level simulation script processes these attempts to propagate the epidemic. We implemented two simulation variants: a finite-population model with susceptible depletion, and an infinite-population model where every infection attempt succeeds. Both are implemented in {\bf \texttt{R}} (version 4.5.0), with performance-critical inner loops reimplemented in C++ via {\bf \texttt{Rcpp}} for speed.

\paragraph{Infection attempt generation.}
For each newly infected individual $i$ (with infection time $t_i$), the infection attempt generator proceeds as follows:
\begin{enumerate}
	\item \textbf{Draw the number of infection attempts.} The number of secondary infection attempts $\chi_i$ is drawn from a $\text{Poisson}(b_i)$ distribution, where $b_i$ is the individual's biological infectiousness. Throughout, we set $b_i \equiv R_0$ for all individuals, so that variation in realized transmission arises solely from infectiousness timing and contact dynamics.
	\item \textbf{Draw the latent period.} A latent period $l_i$ is drawn from a $\text{Gamma}((1-\psi)\alpha, \beta)$ distribution, representing the delay between infection and the onset of transmissibility. When $\psi = 1$, the latent period is zero; when $\psi = 0$, it accounts for the full variance of the generation interval.
	\item \textbf{Draw the timing of each infection attempt.} For each attempt $j = 1, \ldots, \chi_i$, an independent jitter $\epsilon_j$ is drawn from a $\text{Gamma}(\psi\alpha, \beta)$ distribution. The time of infection attempt $j$ relative to the index case's infection time is $\tau_{ij} = l_i + \epsilon_j$. When $\psi = 0$, $\epsilon_j = 0$ and all attempts occur simultaneously at time $l_i$; when $\psi = 1$, the $\epsilon_j$ are independent draws from the full $\text{Gamma}(\alpha, \beta)$ generation interval distribution.
	\item \textbf{Thin for time-varying contacts (if applicable).} When a time-varying contact process $c_i(t)$ is active, each infection attempt at calendar time $t_i + \tau_{ij}$ is accepted with probability $c_i(t_i + \tau_{ij}) / c_{\max}$, where $c_{\max}$ is the supremum of the contact process, using Poisson thinning. For the periodic contact model (Eq.~\ref{eq:periodic_contacts}), all individuals share the same deterministic contact function $c_i(t) = 1 - c_{\text{amp}} \cos(2\pi t / c_{\text{per}})$, and $c_{\max} = 1 + c_{\text{amp}}$. For the stochastic contact model, each individual's piecewise-constant contact trajectory is generated by drawing switching times from a Poisson process with rate $\lambda$ and contact levels from a $\text{Gamma}(\sigma, \sigma)$ distribution; the contact trajectory is generated over an interval $[0, 5\bar{g}]$ to ensure enough contact levels are drawn to accommodate all infection attempts. Thinning is performed against the realized maximum contact level over that individual's trajectory.
	\item \textbf{Thin for detect-and-isolate interventions (if applicable).} When a detect-and-isolate intervention is active, an individual-specific detection time $D_i$ is drawn (relative to peak infectiousness $l_i$) according to the detection mechanism in effect (symptom-based or screening-based). The individual participates with probability $\rho$. If participating, all infection attempts occurring after detection ($\tau_{ij} > t_i + D_i$) are removed with probability $\eta$, representing the effectiveness of post-detection isolation.
\end{enumerate}

\paragraph{Finite-population simulation.}
The finite-population simulation tracks susceptible depletion in a population of size $N$. It proceeds as follows:
\begin{enumerate}
	\item \textbf{Initialization.} A single index case is infected at time $t = 0$. All other $N-1$ individuals are susceptible.
	\item \textbf{Event scheduling.} Each infection attempt is stored as a pending event with its scheduled calendar time $t_i + \tau_{ij}$. Events are maintained in a priority queue, ensuring that the earliest pending event is always processed next.
	\item \textbf{Event processing.} At each step, the earliest event is removed from the queue. A target individual is selected uniformly at random from the population ($N - 1$ individuals excluding the infector). If the target is susceptible, infection occurs: the target is marked as infected, and its own infection attempts are generated and queued. If the target is already infected or recovered, the attempt fails (``wasted'' on a non-susceptible).
	\item \textbf{Termination.} The simulation terminates when the event queue is empty (epidemic extinction) or when a stopping criterion is met (\eg, a specified number of infections or calendar time is reached).
\end{enumerate}

\paragraph{Infinite-population simulation.}
The infinite-population simulation is a special case in which every infection attempt succeeds (no susceptible depletion). This is equivalent to a branching process approximation of the early epidemic and is useful for studying early-epidemic dynamics without the confounding effect of herd immunity. The implementation is identical to the finite-population version except that step 3 is simplified: every queued event produces a new infection, with no target selection or susceptibility check.

\paragraph{Default parameters.}
Unless otherwise stated, simulations used a population size of $N = 10{,}000$, with $n_{\text{sim}} = 5{,}000$ replicate simulations per scenario. Epidemic establishment was defined as reaching $0.05 N = 500$ cumulative infections. To cover a range of infectiousness profile shapes, we considered $\psi \in \{0, 0.2, 0.5, 0.8, 1\}$.

\paragraph{Benchmark pathogens.}
For illustration, we focused on three sets of epidemiological parameters reflecting published estimates of $R_0$ and $g(\tau)$ for influenza, SARS-CoV-2 (omicron), and measles. The generation interval distribution $g(\tau) \sim \text{Gamma}(\alpha, \beta)$ was parameterized using the mean generation time $\bar{g}$ and variance $\text{Var}[g]$ reported in the literature, with $\alpha = \bar{g}^2 / \text{Var}[g]$ and $\beta = \bar{g} / \text{Var}[g]$.

\begin{table}[h]
\centering
\begin{tabular}{lccccc}
\hline
Pathogen & $R_0$ & $\bar{g}$ (days) & $\text{Var}[g]$ (days$^2$) & $\alpha$ & $\beta$ \\
\hline
Influenza & 2 & 3.2 & 4.41 & 2.32 & 0.73 \\
SARS-CoV-2 (omicron) & 6 & 7.05 & 20.8 & 2.39 & 0.34 \\
Measles & 12 & 12.2 & 13.10 & 11.36 & 0.93 \\
\hline
\end{tabular}
\caption{Benchmark pathogen parameters used in simulations.}
\label{tab:benchmark_pathogens}
\end{table}



\clearpage
\subsubsection{Overdispersion in daily infection counts from punctuated infectiousness profiles}
\label{sec:supp_daily_overdispersion}

Here we derive the index of dispersion of the daily infection count attributable to a single infector, show that it is monotonically decreasing in the punctuation parameter $\psi$, and show that the overdispersion aggregates to the population level.

\paragraph{Setup.}
Fix an index case $i$, infected at time $t_i$, who generates $\chi_i \sim \text{Poisson}(R_0)$ total secondary infections. Under the Gamma time-shift model, the generation interval for each secondary infection is $\tau = l_i + \epsilon$, where $l_i \sim \text{Gamma}((1-\psi)\alpha, \beta)$ is a shared latent period and $\epsilon \sim \text{Gamma}(\psi\alpha, \beta)$ is an independent individual jitter. Let $N_i(\tau)$ be the number of person $i$'s secondary infections that fall on day-since-infection $[\tau, \tau+1)$, and define the random landing probability
\begin{equation}
	Q_i(\tau) = \int_\tau^{\tau+1} f_\epsilon(u - l_i) \, du
\end{equation}
where $f_\epsilon$ is the $\text{Gamma}(\psi\alpha, \beta)$ density. The randomness in $Q_i(\tau)$ arises from the draw of $l_i$.

\paragraph{Mean of $Q_i(\tau)$.}
Taking the expectation over $l_i$ and applying Fubini's theorem:
\begin{equation}
	E[Q_i(\tau)] = \int_\tau^{\tau+1} E_{l_i}[f_\epsilon(u - l_i)] \, du = \int_\tau^{\tau+1} (f_l * f_\epsilon)(u) \, du = \int_\tau^{\tau+1} g(u) \, du \equiv q(\tau)
\end{equation}
since $f_l * f_\epsilon = g$ by construction. Thus $E[Q_i(\tau)]$ depends only on the population-level generation interval distribution and is independent of $\psi$.

\paragraph{Index of dispersion of $N_i(\tau)$.}
Conditional on $Q_i(\tau)$, the $\chi_i$ offspring are independently assigned to day $[\tau, \tau+1)$ with probability $Q_i(\tau)$. By the Poisson thinning property (marginalizing over $\chi_i$):
\begin{equation}
	N_i(\tau) \mid Q_i(\tau) \sim \text{Poisson}(R_0 \, Q_i(\tau))
\end{equation}
Applying the law of total variance:
\begin{align}
	\text{Var}[N_i(\tau)] &= E[\text{Var}[N_i(\tau) \mid Q_i(\tau)]] + \text{Var}[E[N_i(\tau) \mid Q_i(\tau)]] \nonumber \\
	&= R_0 \, E[Q_i(\tau)] + R_0^2 \, \text{Var}[Q_i(\tau)]
\end{align}
Since $E[N_i(\tau)] = R_0 \, E[Q_i(\tau)] = R_0 \, q(\tau)$, the index of dispersion is
\begin{equation}
	\text{ID}[N_i(\tau)] = \frac{\text{Var}[N_i(\tau)]}{E[N_i(\tau)]} = 1 + R_0 \, \frac{\text{Var}[Q_i(\tau)]}{E[Q_i(\tau)]}
	\label{eq:ID_individual}
\end{equation}

\paragraph{Extreme cases.}
In the smooth case ($\psi = 1$), $f_\epsilon = g$ and $l_i = 0$ deterministically, so $Q_i(\tau) = q(\tau)$ is constant across individuals. Thus $\text{Var}[Q_i(\tau)] = 0$ and $\text{ID}[N_i(\tau)] = 1$.

In the spike case ($\psi = 0$), $f_\epsilon$ is a point mass at zero. Then $Q_i(\tau) = \mathbf{1}(l_i \in [\tau, \tau+1))$, which is Bernoulli with parameter $q(\tau)$. Thus $\text{Var}[Q_i(\tau)] = q(\tau)(1 - q(\tau))$ and
\begin{equation}
	\text{ID}[N_i(\tau)] = 1 + R_0(1 - q(\tau))
\end{equation}

% \paragraph{Monotonicity of $\text{Var}[Q_i(\tau)]$ in $\psi$.}
% We show that for $\psi_1 < \psi_2$, $\text{Var}[Q_i^{(\psi_1)}(\tau)] \geq \text{Var}[Q_i^{(\psi_2)}(\tau)]$, where the superscript makes the $\psi$-dependence explicit. The proof uses a coupling argument.

% Since Gamma distributions with the same rate parameter are additive, and $(1-\psi_1)\alpha > (1-\psi_2)\alpha$, we can write
% \begin{equation}
% 	L_{\psi_1} \stackrel{d}{=} L_{\psi_2} + \Delta
% \end{equation}
% where $L_{\psi_j} \sim \text{Gamma}((1-\psi_j)\alpha, \beta)$ and $\Delta \sim \text{Gamma}((\psi_2 - \psi_1)\alpha, \beta)$ is independent of $L_{\psi_2}$.

% Under this coupling, we establish a smoothing relationship. Let $Y = \mathbf{1}(\tau_{\text{gen}} \in [\tau, \tau+1))$ where $\tau_{\text{gen}} = L + \epsilon$ is a secondary infection's generation interval. Then $Q_i^{(\psi)}(\tau) = P(\tau_{\text{gen}} \in [\tau, \tau+1) \mid L_\psi) = E[Y \mid L_\psi]$. By the tower property:
% \begin{align}
% 	Q_i^{(\psi_2)}(\tau) &= E[Y \mid L_{\psi_2}] = E\!\left[E[Y \mid L_{\psi_2}, \Delta] \mid L_{\psi_2}\right] \nonumber \\
% 	&= E\!\left[E[Y \mid L_{\psi_1}] \mid L_{\psi_2}\right] = E\!\left[Q_i^{(\psi_1)}(\tau) \mid L_{\psi_2}\right]
% \end{align}
% where the third equality uses the fact that, given $(L_{\psi_2}, \Delta)$, we know $L_{\psi_1} = L_{\psi_2} + \Delta$, and $Y$ depends on $L_{\psi_1}$ and $\epsilon_{\psi_1}$ only, with $\epsilon_{\psi_1}$ independent of $(L_{\psi_2}, \Delta)$.

% Thus $Q_i^{(\psi_2)}(\tau)$ is a conditional expectation of $Q_i^{(\psi_1)}(\tau)$. Applying the law of total variance to $Q_i^{(\psi_1)}(\tau)$, conditioning on $L_{\psi_2}$:
% \begin{equation}
% 	\text{Var}[Q_i^{(\psi_1)}(\tau)] = \text{Var}[Q_i^{(\psi_2)}(\tau)] + E[\text{Var}[Q_i^{(\psi_1)}(\tau) \mid L_{\psi_2}]] \geq \text{Var}[Q_i^{(\psi_2)}(\tau)]
% \end{equation}
% which establishes the claimed monotonicity.

% \paragraph{Aggregation to population-level daily counts.}
% On calendar day $t$, the total number of new infections is $I(t) = \sum_i N_i(t - t_i)$, summing over all previously infected individuals $i$. Conditional on the epidemic history $\{t_i\}$, the $N_i(t - t_i)$ are independent (each infector generates offspring independently). Writing $q_i = q(t - t_i)$ for the population-level landing probability specific to infector $i$:
% \begin{equation}
% 	\text{Var}[I(t) \mid \{t_i\}] = \sum_i \text{Var}[N_i(t - t_i)]
% \end{equation}
% In the smooth case ($\psi = 1$), each $N_i(t-t_i) \sim \text{Poisson}(R_0 q_i)$, so $I(t) \mid \{t_i\} \sim \text{Poisson}(\sum_i R_0 q_i)$ and $\text{ID}[I(t) \mid \{t_i\}] = 1$.

% In the spike case ($\psi = 0$), using $\text{Var}[N_i(t-t_i)] = R_0 q_i + R_0^2 q_i(1-q_i)$:
% \begin{align}
% 	\text{ID}[I(t) \mid \{t_i\}] &= \frac{\sum_i \left[R_0 q_i + R_0^2 q_i(1-q_i)\right]}{\sum_i R_0 q_i} = 1 + R_0\!\left(1 - \frac{\sum_i q_i^2}{\sum_i q_i}\right) = 1 + R_0(1 - \bar{q}_w)
% 	\label{eq:ID_aggregate}
% \end{align}
% where $\bar{q}_w = \sum_i q_i^2 / \sum_i q_i$ is the intensity-weighted mean landing probability. Since $q_i \in [0,1]$, we have $q_i^2 \leq q_i$ for all $i$, so $\bar{q}_w \leq 1$. In a growing epidemic, infectors have a spread of ages since infection, so the $q_i$ values are generally well below 1, making $\bar{q}_w$ small and the overdispersion substantial. The aggregate index of dispersion is bounded: $\text{ID}[I(t) \mid \{t_i\}] \in [1, 1 + R_0]$, with ID $\to 1$ when the generation interval concentrates on a single day and ID $\to 1 + R_0$ when it is spread over many days.

\clearpage

\subsubsection*{Parameter structure}

The punctuated Gamma model has four epidemiological parameters: the basic  reproduction number $R_0$, the early-epidemic growth rate $r$, and the two  parameters of the generation interval distribution, $\alpha$ (shape) and  $\beta$ (rate), which together determine the mean generation time  $\bar{g} = \alpha/\beta$. These four parameters are not all free: the  Euler--Lotka equation
%
\begin{equation}
	1  
	=R_0 \int_{\tau=0}^\infty e^{-r\tau } \frac{\beta^\alpha}{\Gamma(\alpha)} \tau^{\alpha-1} e^{-\beta \tau} d\tau
	= R_0 \left(\frac{\beta}{\beta + r}\right)^\alpha
	\label{eq:euler_lotka}
\end{equation}
%
provides one constraint among them, leaving three free parameters. We  choose to work with $(R_0, \bar{g}, \alpha)$ as the free parameters,  since each has a direct epidemiological interpretation. The rate parameter $\beta$ follows immediately from the  definition of the mean generation interval:
%
\begin{equation}
    \beta = \frac{\alpha}{\bar{g}},
\end{equation}
%
and the early-epidemic growth rate $r$ follows from the Euler--Lotka equation:
%
\begin{equation}
    r = \beta\left(R_0^{1/\alpha} - 1\right) 
      = \frac{\alpha}{\bar{g}}\left(R_0^{1/\alpha} - 1\right).
    \label{eq:r_from_params}
\end{equation}
%
The punctuation  parameter $\psi \in [0, 1]$ is the only remaining free parameter, and  it is the sole quantity that is not identifiable from population-level  epidemic data: since $A(\tau)$ is invariant across $\psi$ by  construction, the parameters $(R_0, \bar{g}, \alpha)$ can in principle  be estimated from the observed epidemic trajectory independently of  $\psi$. The consequences of varying $\psi$ while holding  $(R_0, \bar{g}, \alpha)$ fixed are therefore attributable purely to  changes in the individual-level temporal structure of infectiousness, with no confounding from changes in the population-level transmission  kernel.

\clearpage

\subsubsection*{Derivation of $E[\varsigma]$ and $\text{Var}[\varsigma]$}

{\bf Setup.} Following \citet{morris_computation_2024}, we describe the cumulative infections during the early-epidemic exponential growth phase as a branching process, $Z(t)$. Let $r$ be the epidemic's exponential growth rate; then, as $t \rightarrow \infty$, $e^{-rt} Z(t) \rightarrow W$ almost surely, where $W$ is a random variable that effectively scales the branching process' initial condition, capturing stochastic fluctuations in the early-epidemic process. We denote this convergence as $Z(t) \approx W \cdot e^{rt}$. 

% Here, $r$ is the early-epidemic exponential growth rate and $W$ is a random variable that effectively scales the branching process' initial condition, capturing stochastic fluctuations in the early-epidemic process. 

With a change of variables, we can translate $W$ into a random initial time shift: 
%
\begin{equation}
	Z(t) \approx e^{r (t + \varsigma)} \qquad \text{where} \qquad \varsigma = \frac{\log W - \log E[W]}{r} = \frac{\log W}{r}
\end{equation}
%
when $I_0 = E[W] = 1$ (\ie, when the epidemic is seeded with a single infectious person). Our goal is to determine how $E[\varsigma]$ and $\text{Var}[\varsigma]$ depend on the punctuation parameter $\psi$. To achieve this, we begin by examining the distribution of $W$. 

\noindent {\bf \boldmath Approximating the distribution of $W$.} First, we note that the distribution of $W$ consists of a point mass at 0 with mass $q$ (the probability of epidemic extinction), combined with a smooth density with positive support. The point mass translates into a time shift of $\varsigma = \log(0) / r = -\infty$, corresponding to an epidemic that never establishes. Thus, we restrict our attention to $W | \text{surv}$ and $\varsigma | \text{surv}$, that is, $W$ and $\varsigma$ conditioned on epidemic survival. 

To approximate the distribution of $W|\text{surv}$, we use a $\text{Gamma}(k, \lambda)$ distribution matched to the first two moments. This is a version of the five-moment-matching approach described by \cite{morris_computation_2024}, simplified for analytic tractability. Specifically, we let  
\begin{equation}
	k = \frac{E[W|\text{surv}]^2}{\text{Var}[W|\text{surv}]} \qquad \text{and} \qquad \lambda = \frac{E[W|\text{surv}]}{\text{Var}[W|\text{surv}]}
	\label{eq:k_and_lambda}
\end{equation}
Our task is now to compute $E[W|\text{surv}]$ and $\text{Var}[W|\text{surv}]$. 

\noindent {\bf \boldmath Deriving $E[W|\text{surv}]$.} We note that 
\begin{equation}
E[W] = E[W|\text{surv}] P(\text{surv}) + E[W|\text{extinct}] P(\text{extinct}) = E[W | \text{surv}] \cdot p
\label{eq:p_explanation}
\end{equation}
since $E[W|\text{extinct}] = 0$, and where $p = 1-q$ is the probability of epidemic survival. Thus,
\begin{equation}
\boxed{E[W|\text{surv}] = \frac{1}{p}}
\end{equation}
\begin{center}
{\it [for plugging into expressions for $k$ and $\lambda$, Eq.~\ref{eq:k_and_lambda}]} \\ 
\end{center} 
under our assumption that $E[W] = 1$. 

\noindent {\bf \boldmath Deriving $\text{Var}[W|\text{surv}]$.} The derivation for $\text{Var}[W]$ is somewhat more involved. Recall the Gamma time-shift model, where an infectious person produces $\chi_i \sim \text{Poisson}(R_0)$ secondary infections distributed at times $\tau_j = l_i+ \epsilon_j,\, j\in \{1, \dots, \chi_i\}$ and 
\begin{equation}
l_i \sim \text{Gamma}((1-\psi) \alpha, \beta), \qquad \epsilon_j \sim \text{Gamma}(\psi \alpha, \beta)
\end{equation}
This ensures that a random secondary infection time $\tau$ is distributed according to the generation interval distribution $g(\tau)$, \ie,
\begin{equation}
	\tau \sim \text{Gamma}(\alpha, \beta)
\end{equation}
For notational convenience, we introduce 
\begin{equation}
	E[e^{-c r \tau}] = \int_0^\infty e^{-cr\tau} \frac{\beta^\alpha}{\Gamma(\alpha)} \tau^{\alpha-1} e^{-\beta \tau}\, d\tau
	 = \Bigl(\frac{1}{1 + cz}\Bigr)^\alpha := \rho_c^\alpha
\end{equation}
that is, 
\begin{equation}
\boxed{\rho_c = \frac{1}{1+cz}}
\end{equation}
\begin{center}
{\it [for notational convenience]} \\ 
\end{center}
where $z = r/\beta = r \bar{g} / \alpha = R_0^{1/\alpha} - 1$, a dimensionless quantity that simplifies expressions and allows one to easily plug in different combinations of epidemiological parameters. Here, $\bar{g}$ is the mean of the generation interval distribution $g(\tau)$, and $c$ is a positive integer. The final expression, $z = R_0^{1/\alpha} - 1$, comes from the Euler-Lotka equation \citep{wallinga_how_2007}: 
\begin{equation}
1 = R_0 E[e^{-r \tau}] = R_0 \rho_1^\alpha 
= R_0 \Bigl(\frac{1}{1+z}\Bigr)^\alpha
\label{eq:euler_lotka}
\end{equation}
% \begin{equation}
% 1 = R_0 E[e^{-r \tau}] = R_0 \int_0^\infty e^{-r \tau} \frac{\beta^\alpha}{\Gamma(\alpha)} \tau^{-\alpha-1} e^{-\beta \tau}\, d\tau
% = R_0 \Bigl(\frac{1}{1+z}\Bigr)^\alpha
% \end{equation}
and solving for $z$. 


Next, we introduce the distributional fixed point equation, which is a standard starting point for computing moments of $W$. This equation exploits the self-similarity of the branching process, in which an epidemic initiated by a single individual is, in distribution, a scaled and time-shifted copy of the whole epidemic. At large $t$, the epidemic seeded by the index case $i$ consists of the downstream epidemics caused by their offspring $j$, which have grown approximately to size $Z_j(t) = W_j e^{r(t - \tau_j)}$. That is, 
%
\[
	Z(t) \approx \sum_{j=1}^{\chi_i} W_j e^{rt} e^{-r \tau_j}
\]
Multiplying through by $e^{-rt}$: 
%
\[  W \overset{d}{=} \sum_{j=1}^{\chi_i} e^{-r \tau_j} W_j \]

We now use this expression to calculate $\text{Var}[W]$. Since randomness enters through both $W$ and $\tau$, we first condition on the number of offspring and their infection times, $\tau_\bullet = \{ \chi_i, \tau_1, \tau_2, \dots, \tau_{\chi_i} \}$, using the law of total variance: 
%
\[ \text{Var}[W] = E[\text{Var}\Bigl[\sum_j e^{-r \tau_j} W_j | \tau_\bullet\Bigr]] + \text{Var}[E\Bigl[\sum_j e^{-r \tau_j} W_j | \tau_\bullet\Bigr]] \]
\[
=E\Bigl[ \sum_j e^{-2r\tau_j}  \cdot \text{Var}[W] \Bigr] + 
\text{Var} \Bigl[\sum_j e^{-r \tau_j} \cdot E[W] \Bigr]
\]
\[
= E\Bigl[\sum_j e^{-2 r \tau_j}\Bigr] \cdot \text{Var}[W] + 
(E[W])^2 \cdot \text{Var}\Bigl[\sum_j e^{-r \tau_j}\Bigr]
\]
%
Since $E[W] = I_0 = 1$, we have 
\[
	\text{Var}[W] = \text{Var}[W] E\Bigl[\sum_j e^{-2 r \tau_j}\Bigr] + \text{Var}\Bigl[\sum_j e^{-r \tau_j}\Bigr]
\]
Solving for $\text{Var}[W]$, 
\[
	\text{Var}[W] = \frac{\text{Var}[\sum_j e^{-r \tau_j}]}{1 - E[\sum_j e^{-2 r \tau_j}]}
\]
To simplify the denominator, note that 
\[ 
E\Bigl[\sum_{j=1}^{\chi_i} e^{-2 r \tau_j}\Bigr] = E[\chi_i] E[e^{-2 r \tau_j}]
 = R_0 E[e^{-2 r \tau_j}]
 % = R_0 \Bigl( \frac{\beta}{\beta + 2r}\Bigr )^\alpha 
 = R_0 \rho_2^\alpha 
 \]
Thus, the variance expression simplifies to 
\[
	\text{Var}[W] = \frac{\text{Var}[\sum_j e^{-r \tau_j}]}{1 - R_0 \rho_2^\alpha}
\]
Next, we examine the numerator: 
\[
	\text{Var}[\sum_j e^{-r \tau_j}] = E[(\sum_j e^{-r \tau_j})^2] - E[\sum_j e^{-r \tau_j}]^2 = E[(\sum_j e^{-r \tau_j})^2] - 1
\]
where the last step follows from the Euler-Lotka equation (Eq.~\ref{eq:euler_lotka}). 
% The second term is (R_0 \frac{1}{1+z})
% \[
% 	\text{Var}[\sum_j e^{-r \tau_j}] = E[(\sum_j e^{-r \tau_j})^2] - E[\sum_j e^{-r \tau_j}]^2 = E[(\sum_j e^{- r \tau_j})^2] - 1
% \]
Then, 
\[
	E[(\sum_j e^{- r \tau_j})^2] = E[\sum_j e^{-2 r \tau_j}] + E[\sum_{j \neq k} e^{-r (\tau_j + \tau_k)}]
\]
\[= R_0 \rho_2^\alpha + E[\chi_i (\chi_i - 1)] E[e^{-r \tau_j} e^{-r \tau_k}]\]
\[= R_0 \rho_2^\alpha + R_0^2 E[e^{-r \tau_j} e^{-r \tau_k}]\]
Here, the $\chi_i (\chi_i - 1)$ term enters because there are this many terms in the sum over $j \neq k$. Also, 
\[
E[\chi_i (\chi_i - 1)] = E[\chi_i^2] - E[\chi_i] = \text{Var}[\chi_i] + E[\chi_i]^2 - E[\chi_i] = R_0 + R_0^2 - R_0 = R_0^2
\]
under the assumption that $\chi_i \sim \text{Poisson}(R_0)$. 

Since $\tau_j = l_i + \epsilon_j$, and since $l_i$, $\epsilon_j$, and $\epsilon_k$ are independent, 
\[ E[e^{-r \tau_j} e^{-r \tau_k}]  = E[e^{-2r l_i} e^{-r \epsilon_j} e^{-r \epsilon_k}] = E[e^{-2r l_i}] (E[e^{-r \epsilon}])^2 = \rho_2^{(1-\psi)\alpha} \cdot \rho _1^{2 \psi \alpha}\]
Finally, we can assemble $\text{Var}[W]$: 
\[
	\text{Var}[\sum_j e^{-r \tau_j}] = R_0 \rho_2^\alpha + R_0^2 \rho_2^{(1-\psi)\alpha} \rho_1^{2 \psi \alpha} - 1
\]
and thus, substituting in the $z$ expressions for $\rho_1$ and $\rho_2$ and noting $R_0 = (1 + z)^\alpha$, we obtain
\begin{equation}
	\text{Var}[W] = \frac{
		\Bigl(\frac{1 + z}{1 + 2z}\Bigr)^\alpha + \Bigl(\frac{(1 + z)^2}{1 + 2z}\Bigr)^{(1-\psi) \alpha} - 1
	}{
		1 - \Bigl(\frac{1 + z}{1 + 2z}\Bigr)^\alpha
	} := 
	\frac{\sigma_1^\alpha + \sigma_2^{(1-\psi)\alpha} - 1}{1 - \sigma_1^\alpha}
\end{equation}
where $\sigma_c = \frac{(1 + z)^c}{1 + 2z}$ for positive integer $c$. 

  Last, we note that $E[W^n|\text{surv}] = E[W^n]/p$
  (by the same logic as in Eq.~\ref{eq:p_explanation}), so
  \begin{align}
      \text{Var}[W | \text{surv}]
      &= E[W^2|\text{surv}] - E[W|\text{surv}]^2 = \frac{1}{p} E[W^2] - \frac{1}{p^2} E[W]^2 \\ 
      &= \boxed{\frac{1}{p} (1 + \text{Var}[W]) - \frac{1}{p^2}}
  \end{align}  
% Last, we note that, by the same logic as in Eq.~\ref{eq:p_explanation}, 
% \begin{equation}
% 	\text{Var}[W | \text{surv}] = \frac{1}{p} \cdot \frac{
% 		\Bigl(\frac{1 + z}{1 + 2z}\Bigr)^\alpha + \Bigl(\frac{(1 + z)^2}{1 + 2z}\Bigr)^{(1-\psi) \alpha} - 1
% 	}{
% 		1 - \Bigl(\frac{1 + z}{1 + 2z}\Bigr)^\alpha
% 	}
% \end{equation}
\begin{center}
{\it [for plugging into expressions for $k$ and $\lambda$, Eq.~\ref{eq:k_and_lambda}]}
\end{center}

\noindent {\bf \boldmath Expressions for $k$ and $\lambda$.} Combining the previous derivations and plugging into Eq.~\ref{eq:k_and_lambda} gives 
\begin{equation}
\boxed{k = \frac{1 - \sigma_1^\alpha}{p \sigma_2^{(1 - \psi)\alpha} - 1 + \sigma_1^\alpha},
\qquad
\lambda = p k }
\end{equation}
\begin{center}
	{\it [Parameters of the moment-matched Gamma distribution for W]}
\end{center}

\noindent {\bf \boldmath Deriving the approximate density of $\varsigma$.} Let $f_W(x)$ be the Gamma-distributed approximate density of $W$: 
\begin{equation}
	f_W(x) = \frac{\lambda^k}{\Gamma(k)} x^{k-1} e^{-\lambda x}
\end{equation}
Since $\varsigma = \frac{1}{r} \log W$, we have $W = e^{r \varsigma}$, or, in terms of the arguments of the density functions, $x = e^{r s}$. By the change of variables formula, the density of $\varsigma$ is 
\begin{equation}
f_\varsigma(s) = f_W(e^{rs}) |\frac{dw}{ds}| = 
\frac{\lambda^k}{\Gamma(k)} e^{rs(k-1)} e^{-\lambda e^{rs}} r e^{rs} = \frac{r \lambda^k}{\Gamma(k)} e^{rsk} e^{-\lambda e^{rs}}
\label{eq:varsigma_dist}
\end{equation}
By taking the logarithm, differentiating with respect to $s$, and setting equal to 0, we can determine the mode of the distribution: 
\begin{equation}
	\boxed{\text{Mode}[\varsigma] = \frac{1}{r} \log \Bigl(\frac{k}{\lambda}\Bigr) = \frac{1}{r} \log \frac{1}{p}}
	\label{eq:varsigma_mode}
\end{equation}
\begin{center}
	{\it [Mode of the time shift distribution]}
\end{center}
Thus, the modal time shift depends on the epidemic growth rate $r$ and the establishment probability $p$ (and thus the reproduction number $R_0$), but not on the punctuation parameter $\psi$. 


\noindent {\bf \boldmath Deriving $E[\varsigma]$.} 
It would be possible to derive $E[\varsigma]$ directly from Eq.\ref{eq:varsigma_dist}, but there is a more direct route: since $W \sim \text{Gamma}(k,\lambda)$ and $\varsigma = \frac{1}{r} \log(W)$, it follows that 
\begin{equation}
	E[\varsigma] = \frac{1}{r}(\Psi^{(0)}(k) - \log \lambda)
\end{equation}
where $\Psi^{(0)}(k)$ is the digamma function evaluated at $k$. Plugging in $\lambda = pk$: 
\begin{equation}
\boxed{E[\varsigma] = \frac{\log(1/p)}{r} + \frac{\Psi^{(0)}(k) - \log(k)}{r}}
\end{equation}
\begin{center}
	{\it [Expected value of the time shift distribution]}
\end{center}
The first term is $\text{Mode}[\varsigma]$ (Eq.~\ref{eq:varsigma_mode}), and is independent of the punctuation parameter $\psi$. The second term is a $\psi$-dependent adjustment that is: 
\begin{enumerate}
	\item {\bf  Strictly negative.} {\it Proof:} Let $G \sim \text{Gamma}(k,1)$. Then, $E[G] = k$ and $E[\log G] = \Psi^{(0)}(k)$. By Jensen's inequality, since $\log$ is strictly concave, 
	\[\Psi^{(0)}(k) = E[\log G] < \log(E[G]) = \log k \]
	Thus, $\bigl(\Psi^{(0)}(k) - \log(k)\bigr)/r < 0$. 
	\item {\bf \boldmath Strictly increasing (less negative) in $\psi$.} {\it Proof:} Consider 
	\[\frac{d}{dk} \bigl[ \Psi^{(0)}(k) - \log k \bigr] = 
	\frac{d}{dk} \Biggl[ \int_0^\infty \Bigl(\frac{e^{-t}}{t} - \frac{e^{-kt}}{1 - e^{-t}}\Bigr) dt \Biggr] - \frac{1}{k} 
	= \int_0^\infty \frac{t e^{-k t }}{1 - e^{-t}}\, dt  - \frac{1}{k}
	\]
	Note that $1 - e^{-t} < t$ for $t > 0$, and thus $\frac{t}{1 - e^{-t}} > 1$. So, 
	\[  
	\int_0^\infty \frac{t e^{-kt}}{1 - e^{-t}} > \int_0^\infty e^{-kt} = \frac{1}{k}
	 \]
	 So, 
	 \[\frac{d}{dk} \bigl[ \Psi^{(0)}(k) - \log k \bigr] > 0 \]
	 and thus $(\Psi^{(0)}(k) - \log k)/r$ is increasing in $k$. \\ 
	 Next, we note that $k$ is increasing in $\psi$. The only $\psi$-dependent quantity in $k$ (Eq.\ref{eq:k_and_lambda}) is $\sigma_2^{(1-\psi)\alpha}$, which shows up in the denominator. Since
	 \[\sigma_2 = \frac{(1 + z)^2}{(1 + 2z)} = 1 + \frac{z^2}{1 + 2z} > 1 \quad \text{ for  } z > 0\]
	 the term is maximized when $\psi = 0$ and minimized when $\psi = 1$. Thus, $k$ is minimized when $\psi = 0$ and maximized when $\psi = 1$ --- \ie, $k$ is increasing in $\psi$. 
\end{enumerate}

\noindent Together, these observations demonstrate that the expected time shift sits later than the modal time shift (\ie, the time shift is left-skewed; Observation 1), and the expected time shift moves increasingly later as $a_i(\tau)$ becomes more punctuated (Observation 2). 

\noindent {\bf \boldmath Deriving $\text{Var}[\varsigma]$.} Following again from the properties of $\log W$ where $W$ is $\text{Gamma}(k, \lambda)$-distributed, 
\begin{equation}
	\boxed{\text{Var}[\varsigma] = \frac{1}{r^2} \Psi^{(1)}(k)}
\end{equation}
\begin{center}
	{\it [Variance of the time shift distribution]}
\end{center}
where $\Psi^{(1)}(k)$ is the trigamma function. The trigamma function is strictly decreasing in $k$, since 
\begin{equation}
\frac{d}{dk} \Psi^{(1)}(k) =  \frac{d}{dk }\sum_{n=0}^\infty \frac{1}{(k+n)^2} = 
-2 \sum_{n=0}^\infty \frac{1}{(k + n)^3}
\end{equation}
which is strictly negative when $k > 0$. Thus, $\psi \rightarrow 0$ (punctuated profiles) $\Rightarrow$ $k$ decreasing $\Rightarrow$ $\text{Var}[\varsigma]$ increasing, \ie, more punctuated profiles increase the variance of the time shift. 

\noindent {\bf Summary.} In this section, we derived properties of the random time shift $\varsigma$ by relying on a branching process approximation to the early-epidemic exponential growth phase. When conditioning on epidemic survival, we demonstrated that the time shift becomes increasingly later in expectation and more variable as the individual infectiousness profile $a_i(\tau)$ becomes more punctuated ($\psi \rightarrow 0$).

\clearpage

\subsubsection*{Generality beyond the time-shifted Gamma model}

The preceding derivation used the time-shifted Gamma model and a Gamma-distributed generation interval. Here, we show that the core result --- punctuated profiles increase the variance of early-epidemic timing --- holds for \emph{any} non-degenerate generation interval distribution (\ie, any distribution with positive variance, as opposed to a degenerate distribution concentrated on a single point). The argument rests only on the additive splitting $\tau_j = l_i + \varepsilon_j$ (shared shift plus independent jitter) and the Cauchy-Schwarz inequality.

\noindent {\bf Setup.} Let $g$ be an arbitrary generation interval distribution with Laplace transform $M_g(s) = E[e^{-s\tau}]$, so the Euler-Lotka equation gives $R_0 M_g(r) = 1$. Each index case's offspring share a common latent shift $l_i$ and receive independent jitter $\varepsilon_j$, so $\tau_j = l_i + \varepsilon_j$ with the marginal distribution of $\tau_j$ equal to $g$. The distributions of $l$ and $\varepsilon$ are arbitrary, subject only to $M_l \cdot M_\varepsilon = M_g$.

\noindent {\bf \boldmath Model-free $\text{Var}[W]$ formula.} The law-of-total-variance derivation proceeds identically to the Gamma case. The only splitting-dependent step is the cross-term for $j \neq k$:
\[
E[e^{-r \tau_j} e^{-r \tau_k}] = E[e^{-2rl}] (E[e^{-r\varepsilon}])^2 = M_l(2r) \cdot M_\varepsilon(r)^2
\]
Using $M_l(2r) = M_g(2r) / M_\varepsilon(2r)$, this becomes $M_g(2r) \cdot Q$, where
\begin{equation}
	Q := \frac{M_\varepsilon(r)^2}{M_\varepsilon(2r)}
	\label{eq:Q_def}
\end{equation}
Assembling the full variance formula (with $A := R_0 M_g(2r)$):
\begin{equation}
	\text{Var}[W] = \frac{A(1 + R_0 Q) - 1}{1 - A}
	\label{eq:var_W_general}
\end{equation}
Here, $A$ depends only on $g$ and not on the splitting, and $Q$ is the \emph{only} place the splitting enters. Since Eq.~\ref{eq:var_W_general} is affine in $Q$ with positive slope $R_0 A / (1-A)$, $\text{Var}[W]$ is strictly increasing in $Q$.

\noindent {\bf Three properties of $Q$.}

\begin{enumerate}
	\item $Q \leq 1$, with equality if and only if $\varepsilon$ is degenerate (maximum punctuation).

	{\it Proof:} By the Cauchy-Schwarz inequality applied to $X = e^{-r\varepsilon}$:
	\[M_\varepsilon(2r) = E[X^2] \geq (E[X])^2 = M_\varepsilon(r)^2\]
	so $Q = M_\varepsilon(r)^2 / M_\varepsilon(2r) \leq 1$. Equality holds if and only if $X$ is almost surely constant, \ie, $\varepsilon$ is degenerate. $\blacksquare$

	\item $Q$ is multiplicative under convolution: if $\varepsilon_2 \overset{d}{=} \varepsilon_1 + \delta$ with $\delta$ independent, then $Q_2 = Q_1 \cdot Q_\delta \leq Q_1$, with strict inequality if $\delta$ is non-degenerate.

	{\it Proof:} $M_{\varepsilon_2}(s) = M_{\varepsilon_1}(s) \cdot M_\delta(s)$, so $Q_2 = Q_1 \cdot M_\delta(r)^2 / M_\delta(2r) = Q_1 \cdot Q_\delta$, and $Q_\delta \leq 1$ by Property 1. $\blacksquare$

	\item The spike-vs-smooth gap in $\text{Var}[W]$ is strictly positive for any non-degenerate $g$.

	{\it Proof:} For the spike ($\varepsilon = 0$), $Q_\text{spike} = 1$. For the smooth ($\varepsilon = \tau$, $l = 0$), $Q_\text{smooth} = M_g(r)^2 / M_g(2r) = 1 / (R_0^2 M_g(2r)) = 1/(R_0 A)$. The gap is
	\[
	\text{Var}[W]_\text{spike} - \text{Var}[W]_\text{smooth} = \frac{R_0 A (1 - Q_\text{smooth})}{1 - A} = \frac{R_0 A - 1}{1 - A}
	\]
	Since $R_0 A = R_0^2 M_g(2r) \geq R_0^2 M_g(r)^2 = 1$ by the Cauchy-Schwarz inequality applied to $X = e^{-r\tau}$, with strict inequality for non-degenerate $g$, the gap is strictly positive. $\blacksquare$
\end{enumerate}

\noindent {\bf Monotonicity for parameterized families.} For a one-parameter family of splittings $\varepsilon_\psi$ with $\varepsilon_0 = 0$ (spike) and $\varepsilon_1 = \tau$ (smooth), if the family is ordered by convolution --- meaning $\varepsilon_{\psi_2} \overset{d}{=} \varepsilon_{\psi_1} + \delta$ for some independent $\delta \geq 0$ whenever $\psi_2 > \psi_1$ --- then $Q(\psi)$ is strictly decreasing and $\text{Var}[W](\psi)$ is strictly decreasing in $\psi$. The time-shifted Gamma model satisfies this ordering, since $\text{Gamma}(\psi_2 \alpha, \beta) \overset{d}{=} \text{Gamma}(\psi_1 \alpha, \beta) + \text{Gamma}((\psi_2 - \psi_1) \alpha, \beta)$, but any infinitely divisible generation interval distribution admits such a family.

\noindent {\bf Implications for the full chain.} Combined with the model-free conditioning step ($\text{Var}[W|\text{surv}] = (1 + \text{Var}[W])/p - 1/p^2$, which increases in $\text{Var}[W]$) and the Gamma-approximation results above, the full chain holds for general $g$:
\begin{center}
\begin{tabular}{ccccccc}
	More punctuation & $\Rightarrow$ & higher $Q$ & $\Rightarrow$ & higher $\text{Var}[W]$ & $\Rightarrow$ & lower $k$ \\
	& & {\scriptsize (Cauchy-Schwarz)} & & {\scriptsize (Eq.~\ref{eq:var_W_general})} & & {\scriptsize (Eq.~\ref{eq:k_and_lambda})} \\[6pt]
	& $\Rightarrow$ & more negative $E[\varsigma]$ & and & higher $\text{Var}[\varsigma]$ \\
	& & {\scriptsize (digamma monotonicity)} & & {\scriptsize (trigamma monotonicity)} \\
\end{tabular}
\end{center}
The first two links are model-free: the only assumption is that $g$ is non-degenerate. The final link --- from $\text{Var}[W]$ to properties of $\varsigma = (\log W)/r$ --- requires a distributional assumption, because moments of $W$ do not uniquely determine moments of $\log W$. The Gamma moment-matching approximation serves this role: given $E[W|\text{surv}]$ and $\text{Var}[W|\text{surv}]$, it provides a complete distribution from which $E[\varsigma]$ and $\text{Var}[\varsigma]$ can be computed in closed form via the digamma and trigamma functions. Without this (or an analogous) approximation, the qualitative direction is still supported by Jensen's inequality --- which guarantees $E[\varsigma|\text{surv}] < \log(1/p)/r$ for any distribution of $W|\text{surv}$, with the gap increasing when the distribution becomes more dispersed in the sense of convex order --- but the specific quantitative expressions would not apply.


\subsubsection*{Calculating overdispersion}

{\bf \boldmath Deriving an expression for the overdispersion parameter $k$.} We begin with the standard assumption that the number of offspring $Z_i$ produced by an index case $i$ is Poisson-distributed: 
\begin{equation}
	Z_i \sim \text{Poisson}(\nu_i)
\end{equation}
When $\nu_i$ is uniform across people ($\nu_i \equiv R_0$), then the $Z_i$ are truly Poisson-distributed. Overdispersion occurs when $\nu_i$ varies across individuals. To capture this overdispersion, the standard approach is to describe $Z_i$ using a Negative Binomial distribution: 
\begin{equation}
	Z_i \sim \text{NegBin}(R_0, k)
\end{equation}
which has expectation $R_0$ and variance $R_0 (1 + R_0 / k)$. Thus, $k$ captures overdispersion, with $k \rightarrow \infty$ yielding Poisson-distributed $Z_i$ and $k \rightarrow 0$ yielding increasingly heavier Negative Binomial tails (\ie, more overdispersion). The Negative Binomial expression is exact when $\nu_i$ are Gamma-distributed; otherwise, it is an approximation to the distribution of $Z_i$. Either way, the parameter $k$ can be obtained {\it via} moment matching by re-arranging the variance expression: 
\begin{equation}
	k = \frac{R_0^2}{\text{Var}[Z_i] - R_0}
\end{equation}
Furthermore, since
\begin{equation}
	\text{Var}[Z_i] = \text{Var}[E[Z_i | \nu_i]] + E[\text{Var}[Z_i | \nu_i]] = \text{Var}[\nu_i] + E[\nu_i] = \text{Var}[\nu_i] + R_0
\end{equation}
the expression for $k$ simplifies to 
\begin{equation}
	k = \frac{R_0^2}{\text{Var}[\nu_i]}
	\label{eq:k_expression}
\end{equation}

Next, if we assume everyone's inherent biological infectiousness is the same ($b_i \equiv R_0$), we have 
\begin{equation}
	\nu_i = R_0 \int_0^\infty g_i(\tau) c_i(t_i + \tau) d\tau := R_0 \tilde{c}_i
\end{equation}
Plugging in to \ref{eq:k_expression} gives 
\begin{equation}
	\boxed{k = \frac{1}{\text{Var}[\tilde{c}_i]}}
\end{equation}
Thus, the overdispersion $k$ depends only on the variance of $\tilde{c}_i = \int_0^\infty g_i(\tau) c_i(t_i + \tau) d\tau$. 

\noindent {\bf \boldmath Deriving an expression for $\text{Var}[\tilde{c}_i]$.} The random variable $\tilde{c}_i$ inherits randomness from both $g_i(\tau)$ and $c_i(t_i+\tau)$. To manage this, we again apply the Law of Total Variance: 
\begin{align}
	\text{Var}[\tilde{c}_i] &= \text{Var}[E[\tilde{c}_i | g_i ]] + E[\text{Var}[\tilde{c}_i | g_i ]] \\ 
	&= \text{Var}[1] + E[\text{Var}[\tilde{c}_i | g_i ]] \\ 
	&= E[\text{Var}[\tilde{c}_i | g_i ]]
\end{align}
Thus, we are interested in 
\begin{equation}
\tilde{c}_i | g_i = \int_0^\infty g_i(\tau) c_i(t_i+\tau) d\tau 
\end{equation}
where $g_i$ is known (\ie, conditioned upon). This has the form of a classic problem in signal processing: $c_i$ is a wide-sense stationary stochastic process (an input signal); $g_i$ is an impulse response; and $\tilde{c}_i | g_i$ is the output signal evaluated at time $t_i$. Together, this comprises a linear system. Following standard theory, we can express $\text{Var}[\tilde{c}_i | g_i]$ in terms of the autocorrelation $R(\tau)$ of the signal $\tilde{c}_i | g_i$ at time lag $\tau$. Specifically, 
\begin{align}
	\text{Var}[\tilde{c}_i | g_i] &= E[(\tilde{c}_i | g_i)^2]  - E[\tilde{c}_i | g_i]^2  \\ 
	&= 	R_{\tilde{c}_i | g_i}(0) - E[\tilde{c}_i | g_i]^2 \\
	&= 	R_{\tilde{c}_i | g_i}(0) - 1 
\end{align}

\noindent {\bf \boldmath Deriving an expression for $R_{\tilde{c}_i | g_i}(0)$.} Deriving an expression for the second moment of the output process, $R_{\tilde{c}_i | g_i}(0)$, is easiest in frequency space. We can introduce $S_\bullet$ as the Fourier transform of $R_\bullet$ and $\hat{g}_i$ as the Fourier transform of $g_i$, \ie 
\begin{equation}
S_\bullet(\omega) = \int_{-\infty}^{\infty} e^{-i \omega \tau} R_\bullet(\tau) d\tau 
\qquad \text{ and } \qquad 
\hat{g}_i(\omega) = \int_{-\infty}^{\infty} e^{-i \omega \tau} g_i(\tau) d\tau
\end{equation}
By the inverse Fourier transform, 
\begin{equation}
R_\bullet(\tau) = \frac{1}{2\pi} \int_{-\infty}^{\infty} S_\bullet(\omega) e^{i \omega \tau} d\omega
\end{equation}
and thus 
\begin{equation}
R_\bullet(0) = \frac{1}{2\pi} \int_{-\infty}^{\infty} S_\bullet(\omega) d\omega
\end{equation}
We thus seek an expression for $S_{\tilde{c}_i | g_i}(\omega)$, which we can obtain from the fundamental theorem for linear systems: 
\begin{equation}
	S_{\tilde{c}_i | g_i}(\omega) = |\hat{g}_i(\omega)|^2 S_{c_i}(\omega)
\end{equation}

\noindent {\bf \boldmath Deriving $\hat{g}_i(\omega).$} We begin by examining $\hat{g}_i(\omega)$. Under the Gamma time-shift model, $g_i(\tau) = f_\epsilon(\tau - l_i)$, where $f_\epsilon$ is the $\text{Gamma}(\psi \alpha, \beta)$ density and $l_i \sim \text{Gamma}((1-\psi) \alpha, \beta)$. Taking the Fourier transform, 
\begin{align}
	\hat{g}_i(\omega) &= \int_{-\infty}^{\infty} g_i(\tau) e^{-i \omega \tau} d\tau \\ 
	&=	\int_{l_i}^{\infty} f_\epsilon(\tau - l_i) e^{-i \omega \tau}  d\tau\\ 
	&= \int_0^\infty f_\epsilon(\tau) e^{-i \omega (\tau + l_i)} d\tau \\
	&=  e^{-i \omega l_i} \int_0^\infty f_\epsilon(\tau) e^{-i \omega \tau} d\tau \\ 
	&= e^{-i \omega l_i} \int_0^\infty \frac{\beta^{\psi \alpha}}{\Gamma(\psi \alpha)} \tau^{\psi \alpha - 1} e^{-\beta \tau} e^{-i\omega\tau} d\tau \\
&= e^{-i \omega l_i} \Bigl(  \frac{\beta}{\beta + i \omega}\Bigr)^{\psi \alpha}
\end{align}
Taking the squared modulus gives 
\begin{equation}
 |\hat{g}_i(\omega)|^2 = \Bigl(1 + \frac{\omega^2}{\beta^2}\Bigr)^{-\psi \alpha}
\end{equation}
This is the power transfer function; it is constant across individuals and does not depend on $l_i$. 

\noindent {\bf \boldmath Deriving $S_{c_i}(\omega)$ for the periodic contact model.} For the periodic contact model
\begin{equation}
c_i(t) = 1 - c_{amp} \cos(\omega_0 t)
\end{equation}
the autocorrelation is 
\begin{align}
R_{c_i}(\tau)  &= E[(1 - c_{amp} \cos(\omega_0 t)) (1 - c_{amp} \cos(\omega_0 (t + \tau)))] \\ 
&= E[1] - c_{amp}E[\cos(\omega_0 t)] - c_{amp}E[\cos(\omega_0 (t + \tau))] + c_{amp}^2 E[\cos(\omega_0 t) \cos(\omega_0 (t + \tau))] \\ 
&= 1 + \frac{c_{amp}^2}{2} \cos(\omega_0 \tau) 
\end{align}
Taking the Fourier transform, 
\begin{align}
	S_{c_i}(\omega) &= \int_{-\infty}^{\infty} \Bigl[ 1 + \frac{c_{amp}^2}{2} \cos(\omega_0\tau) \Bigr] e^{-i\omega\tau}d\tau \\ 
	&= 2 \pi \delta(\omega) + \frac{\pi c_{amp}^2}{2}[\delta(\omega - \omega_0) + \delta(\omega + \omega_0)] 
\end{align}

\noindent {\bf \boldmath Constructing $k$ for the periodic contact model.} Putting everything together, we have: 
\begin{align}
\text{Var}[\tilde{c}_i | g_i] &= R_{\tilde{c}_i | g_i}(0) - 1 \\
&= \frac{1}{2\pi} \int_{-\infty}^{\infty} S_{\tilde{c}_i | g_i}(\omega) d\omega  - 1 \\
&= \frac{1}{2\pi} \int_{-\infty}^\infty \Bigl(1 + \frac{\omega^2}{\beta^2}\Bigr)^{-\psi \alpha} 
\Bigl[ 2 \pi \delta(\omega) + \frac{\pi c_{amp}^2}{2} [\delta(\omega - \omega_0) + \delta(\omega + \omega_0)] \Bigr] d \omega  - 1\\ 
&= 1 + \frac{c_{amp}^2}{2} \Bigl(1 + \frac{\omega_0^2}{\beta^2}\Bigr)^{-\psi \alpha} - 1 \\
&= \frac{c_{amp}^2}{2} \Bigl(1 + \frac{\omega_0^2}{\beta^2}\Bigr)^{-\psi \alpha}
\end{align}
Thus, 
\begin{equation}
	\boxed{k = \frac{2}{c_{amp}^2} \Bigl(1 + \frac{\omega_0^2}{\beta^2}\Bigr)^{\psi \alpha}}
\end{equation}

\noindent {\bf \boldmath Deriving $S_{c_i}(\omega)$ for the Gamma-Poisson contact model.} Let $X \sim \text{Gamma}(k_c, k_c)$ be a contact-level draw from the Gamma-Poisson contact model. Then, the autocorrelation 
\begin{equation}
	R_{c_i}(\tau) = E[c_i(t) c_i(t + \tau)]
\end{equation}
is equal to $E[X^2] = \text{Var}[X] + E[X]^2 = \frac{1}{k_c} + 1$ with probability $e^{-\lambda |\tau|}$ (\ie, if the contact levels have not switched between times $t$ and $t + \tau$), and is equal to $E[X]^2 = 1$ otherwise. That is: 
\begin{align}
	R_{c_i}(\tau) &= (\frac{1}{k_c} +1) e^{-\lambda |\tau|} + (1)(1 - e^{-\lambda |\tau|}) \\ 
	&= \frac{1}{k_c} e^{-\lambda |\tau|} + 1
\end{align}
Thus: 
\begin{align}
S_{c_i}(\omega) &= \int_{-\infty}^{\infty} [\frac{1}{k_c}e^{-\lambda |\tau| } + 1] e^{-i \omega \tau} d\tau \\ 
&= \frac{1}{k_c} \int_{-\infty}^{0} e^{ \lambda \tau} e^{- i \omega \tau} d\tau + 
\frac{1}{k_c} \int_{0}^{\infty} e^{- \lambda \tau } e^{- i \omega \tau} d\tau + 2\pi \delta(\omega) \\ 
&= \frac{1}{k_c} \Bigl[ \frac{1}{\lambda - i \omega} + \frac{1}{\lambda + i \omega} \Bigr]  + 2\pi \delta(\omega) \\
&= \frac{2 \lambda}{k_c (\lambda^2 + \omega^2)} + 2\pi \delta(\omega) 
\end{align}

\noindent {\bf \boldmath Constructing $k$ for the Gamma-Poisson contact model.} Putting everything together, we have: 
\begin{align}
\text{Var}[\tilde{c}_i | g_i] &= R_{\tilde{c}_i | g_i}(0) - 1 \\
&= \frac{1}{2\pi} \int_{-\infty}^{\infty} 
\Bigl[\frac{2 \lambda}{k_c (\lambda^2 + \omega^2)} + 2\pi \delta(\omega)\Bigr]
\Bigl( 1 + \frac{\omega^2}{\beta^2}\Bigr )^{-\psi \alpha} d\omega - 1 \\
&= \frac{1}{2\pi} \int_{-\infty}^{\infty} 
\Bigl[\frac{2 \lambda}{k_c (\lambda^2 + \omega^2)} \Bigr] \Bigl( 1 + \frac{\omega^2}{\beta^2}\Bigr )^{-\psi \alpha} d\omega \\
&= \frac{1}{k_c}\int_{-\infty}^{\infty} \frac{\lambda}{\pi (\lambda^2 + \omega^2)} \Bigl( 1 + \frac{\omega^2}{\beta^2}\Bigr )^{-\psi \alpha} d\omega
\end{align}
and $k$ is the reciprocal of this integral. The integral can be solved in closed form using special functions, but its final form is not especially illuminating; in practice, we evaluate this expression by numerically solving the integral. 

\clearpage

\subsection{Supplementary Figures}